\documentclass[a4paper,12pt]{article}

\usepackage[utf8]{inputenc}
\usepackage[T2A]{fontenc}
\usepackage[russian]{babel}
\usepackage{tikz}
\usetikzlibrary{patterns}
\usepackage{geometry}
\usepackage{xcolor}
\usepackage{amssymb}
\usepackage{paracol}
\usepackage{wrapfig}

\geometry{margin=1.5cm, left=10mm, right=10mm}
\pagestyle{empty}

\newcommand{\seg}[3][black]{%
\begin{tikzpicture}[baseline=-0.5ex]
  \draw[#1, line width=2.5pt] (0,0) -- (1.2,0);
  \node[above, font=\scriptsize] at (0,0.05) {#2};
  \node[above, font=\scriptsize] at (1.2,0.05) {#3};
\end{tikzpicture}%
}

\newcommand{\segsm}[3][black]{%
\begin{tikzpicture}[baseline=-0.5ex]
  \draw[#1, line width=2.5pt] (0,0) -- (0.8,0);
  \node[above, font=\scriptsize] at (0,0.05) {#2};
  \node[above, font=\scriptsize] at (0.8,0.05) {#3};
\end{tikzpicture}%
}

\newcommand{\mixsegpow}[6][black]{%
\begin{tikzpicture}[baseline=0.2ex]
  \draw[#1, line width=2.5pt] (0,0) -- (#2,0);
  \draw[#1, line width=2.5pt, dash pattern=on 3pt off 3pt] (#2,0) -- ({#2+#3},0);
  \node[above, font=\scriptsize] at (0,0.05) {#4};
  \node[above, font=\scriptsize] at ({#2+#3},0.05) {#5};
\end{tikzpicture}^{#6}%
}

\newcommand{\mixseg}[5][black]{%
\begin{tikzpicture}[baseline=-0.5ex]
  \draw[#1, line width=2.5pt] (0,0) -- (#2,0);
  \draw[#1, line width=2.5pt, dash pattern=on 3pt off 3pt] (#2,0) -- ({#2+#3},0);
  \node[above, font=\scriptsize] at (0,0.05) {#4};
  \node[above, font=\scriptsize] at ({#2+#3},0.05) {#5};
\end{tikzpicture}%
}

\newcommand{\dsegsm}[3][black]{%
\begin{tikzpicture}[baseline=-0.5ex]
  \draw[#1, line width=2.5pt, dash pattern=on 3pt off 3pt] (0,0) -- (0.8,0);
  \node[above, font=\scriptsize] at (0,0.05) {#2};
  \node[above, font=\scriptsize] at (0.8,0.05) {#3};
\end{tikzpicture}%
}

\begin{document}

\noindent
\textbf{КНИГА II ПРЕДЛ. XIII. ТЕОРЕМА} \hfill 89

\vspace{0.5cm}

\columnratio{0.65}
\begin{paracol}{2}

\begin{wrapfigure}{l}{0.10\textwidth}
    \includegraphics[width=0.10\textwidth]{image.png}
\end{wrapfigure}

\noindent
\textit{о всяком треугольнике квадрат стороны, стягивающей острый угол, меньше суммы квадратов сторон, содержащих этот угол, на дважды прямоугольник, заключенный между любой из этих сторон и отрезком, отсекаемым перпендикуляром из противоположного угла от этого отрезка или от продленного отрезка.}

\vspace{0.3cm}

\textit{Первый случай.}

$\seg[blue]{C}{A}^2 < \mixsegpow[black]{0.7}{0.5}{B}{C}{2} + \seg[red]{A}{B}^2$ на $2 \cdot \mixseg[black]{0.7}{0.5}{B}{C} \cdot \segsm[black]{B}{D}$.

\vspace{0.3cm}

\textit{Второй случай.}

$\seg[blue]{E}{G}^2 < \seg[red]{E}{F}^2 + \seg[black]{F}{G}^2$ на $2 \cdot \seg[black]{F}{G} \cdot \dsegsm[black]{F}{H}$.

\vspace{0.5cm}

\textit{Предположим, перпендикуляр падает внутри треугольника, тогда (пр. II.7)}

\vspace{0.2cm}

$\mixsegpow[black]{0.7}{0.5}{B}{C}{2} + \segsm[black]{B}{D}^2 = 2 \cdot \mixseg[black]{0.7}{0.5}{B}{C} \cdot \segsm[black]{B}{D} + \dsegsm[black]{D}{C}^2$,

\vspace{0.2cm}

\textit{к каждой добавим $\seg[orange]{A}{D}^2$, тогда}

\vspace{0.2cm}

$\mixsegpow[black]{0.7}{0.5}{B}{C}{2} + \segsm[black]{B}{D}^2 + \seg[orange]{A}{D}^2 = 2 \cdot \mixseg[black]{0.7}{0.5}{B}{C} \cdot \segsm[black]{B}{D} + \dsegsm[black]{D}{C}^2 + \seg[orange]{A}{D}^2$

\vspace{0.2cm}

$\therefore$ (пр. I.47) $\mixsegpow[black]{0.7}{0.5}{B}{C}{2} + \seg[red]{A}{B}^2 = 2 \cdot \mixseg[black]{0.7}{0.5}{B}{C} \cdot \segsm[black]{B}{D} + \seg[blue]{C}{A}^2$,

\vspace{0.2cm}

и $\therefore$ $\seg[blue]{C}{A}^2 < \mixsegpow[black]{0.7}{0.5}{B}{C}{2} + \seg[red]{A}{B}^2$ на $2 \cdot \mixseg[black]{0.7}{0.5}{B}{C} \cdot \dsegsm[black]{D}{C}$.

\vspace{0.5cm}

\textit{Теперь предположим, что перпендикуляр падает вовне треугольника, тогда (пр. II.7)}

\vspace{0.2cm}

$\mixsegpow[black]{0.7}{0.5}{F}{H}{2} + \seg[black]{F}{G}^2 = 2 \cdot \mixseg[black]{0.7}{0.5}{F}{H} \cdot \segsm[black]{F}{G} + \dsegsm[black]{G}{H}^2$,

\vspace{0.2cm}

\textit{к каждой добавим $\seg[orange]{H}{E}^2$, тогда}

\vspace{0.2cm}

$\mixsegpow[black]{0.7}{0.5}{F}{H}{2} + \seg[black]{F}{G}^2 + \seg[orange]{H}{E}^2 = 2 \cdot \mixseg[black]{0.7}{0.5}{F}{H} \cdot \segsm[black]{F}{G} + \dsegsm[black]{G}{H}^2 + \seg[orange]{H}{E}^2$

\vspace{0.2cm}

$\therefore$ (пр. I.47) $\seg[red]{E}{F}^2 + \seg[black]{F}{G}^2 = 2 \cdot \mixseg[black]{0.7}{0.5}{F}{H} \cdot \segsm[black]{F}{G} + \seg[blue]{E}{G}^2$,

\vspace{0.2cm}

$\therefore$ $\seg[blue]{E}{G}^2 < \seg[red]{E}{F}^2 + \seg[black]{F}{G}^2$ на $2 \cdot \mixseg[black]{0.7}{0.5}{F}{H} \cdot \segsm[black]{F}{G}$.

\vspace{0.5cm}

\hfill ч. т. д.

\switchcolumn

\vspace{0.5cm}

\centering
\textit{Первый случай}

\vspace{0.3cm}

\begin{tikzpicture}[scale=1.2]
  \coordinate (A) at (1.5,3.5);
  \coordinate (B) at (0,0);
  \coordinate (C) at (3.5,0);
  \coordinate (D) at (1.5,0);
  
  \draw[red, line width=3pt] (B) -- (A);
  \draw[blue, line width=3pt] (A) -- (C);
  \draw[black, line width=3pt] (B) -- (D);
  \draw[black, line width=3pt, dash pattern=on 3pt off 3pt] (D) -- (C);
  \draw[orange, line width=3pt] (A) -- (D);
  
  \node[above] at (A) {A};
  \node[below left] at (B) {B};
  \node[below right] at (C) {C};
  \node[below] at (D) {D};
\end{tikzpicture}

\vspace{1.5cm}

\textit{Второй случай}

\vspace{0.3cm}

\begin{tikzpicture}[scale=1.2]
  \coordinate (E) at (3.5,3.5);
  \coordinate (F) at (0,0);
  \coordinate (H) at (3.5,0);
  \coordinate (G) at (1.5,0);
  
  \draw[red, line width=3pt] (F) -- (E);
  \draw[orange, line width=3pt] (E) -- (H);
  \draw[black, line width=3pt] (F) -- (G);
  \draw[black, line width=3pt, dash pattern=on 3pt off 3pt] (G) -- (H);
  \draw[blue, line width=3pt] (E) -- (G);
  
  \node[above right] at (E) {E};
  \node[below left] at (F) {F};
  \node[below right] at (H) {H};
  \node[below] at (G) {G};
\end{tikzpicture}

\end{paracol}

\end{document}